
\section{Preliminary}
\subsection{Optimal Transport}
Optimal Transport (OT)~\cite{solomon2018optimal} provides a principled framework for measuring the
discrepancy between two distributions while explicitly accounting for the
underlying geometry. Assume two uniform discrete distributions $\mathbf{a}$ and $\mathbf{b}$, supported on $T_m$ and $T_t$ elements, respectively. We denote the set of admissible transport plans by
\begin{equation}
\label{eq:P_constraint}
U(\mathbf{a},\mathbf{b})=
\left\{
\mathbf{P} \in \mathbb{R}_+^{T_m \times T_t}
\;\middle|\;
\mathbf{P}\mathbf{1}_{T_t}=\mathbf{a},\;
\mathbf{P}^\top \mathbf{1}_{T_m}=\mathbf{b}
\right\},
\end{equation}
where
\begin{equation}
\mathbf{a}=\frac{1}{T_m}\mathbf{1}_{T_m},
\qquad
\mathbf{b}=\frac{1}{T_t}\mathbf{1}_{T_t}.
\end{equation}

The marginal constraints ensure that $\mathbf{P}$ redistributes the mass of $\mathbf{a}$ onto $\mathbf{b}$.
Given a cost matrix $\mathbf{C}\in\mathbb{R}^{T_m\times T_t}$ measuring the cost of
transporting unit mass between elements, the Kantorovich formulation of OT is
defined as
\begin{equation}
\label{eq:ot_kantorovich}
\min_{\mathbf{P}\in{U}(\mathbf{a},\mathbf{b})}
\ \langle \mathbf{P},\mathbf{C}\rangle
\end{equation}
where $\langle \mathbf{P},\mathbf{C}\rangle=\sum_{i,j}\mathbf{P}_{ij}\mathbf{C}_{ij}$.

In contrast to pointwise matching, OT naturally supports many-to-many
correspondence and enforces global consistency via marginal constraints.
These properties make OT particularly suitable for modeling fine-grained
alignment under weak supervision, such as motion--text correspondence.

\subsection{Entropic-Regularized Optimal Transport and Sinkhorn Algorithm}
\label{sec:entro}

Directly solving \cref{eq:ot_kantorovich} requires linear programming, which is
computationally expensive and unsuitable for end-to-end learning.
To address this issue, entropic regularization augments the OT objective with an
entropy term:
\begin{equation}
\label{eq:entropic_ot}
\begin{aligned}
\mathbf{P}^* = \arg\min_{\mathbf{P}\in{U}(\mathbf{a},\mathbf{b})}
\ & \langle \mathbf{P},\mathbf{C}\rangle
\;+\;
\varepsilon
\sum_{i,j}{P}_{ij}\big(\log{P}_{ij}-1\big)
\end{aligned}
\end{equation}
where $\varepsilon>0$ controls the strength of the regularization. The entropic regularization renders the problem strictly convex and enables efficient optimization via the Sinkhorn algorithm.
%As a result, entropic OT can be seamlessly integrated as a loss function or alignment module in deep learning models.
